\chapter{Always be somewhat suspect (Magic and Miracles Are Real)}
The fourth episode of Madoka Magica begins the second arc of the show, but the main foci of that arc—Sayaka as a magical girl and the new character Kyoko—do not appear until the end of the episode. The majority of the episode serves instead as a coda to the first arc and an introduction to the second.

In addition, a case can be made that this is the first "true" episode of the series; that is, that this is the first episode which openly presents its dominant aesthetic, as opposed to attempting to pretend to be a more typical magical girl show in the first three episodes. It is thus also an introduction to the series itself, as it presents for the first time one of the most important themes of the series: depression.

There's a quote from series writer Gen Urobuchi that is often cited in reference to Madoka Magica, though the quote itself is from the afterword to the first volume of his earlier work Fate/zero. In part, the quote reads, "The truth is, I haven't always been this way. I have often written pieces that didn't have a perfect ending, but by the last chapter the protagonist would still possess a belief that 'Although there will be many hardships to come, I still have to hold on.' But ever since I don't know when, I can no longer write works like this." \cite{ref37} He goes on to talk about the inevitability of entropy and failure, in a way suggestive of some sort of depression—or at least despair.

Now, it is always dangerous to make guesses at an author's mental state or opinions based on their work. The implied author (the answer to the question "what kind of person would write this work" \cite{ref38}) is inevitably different from the actual author, who like all real people (and unlike fictional people such as either characters or implied authors) has an entirely unknowable subjective internality that cannot be perceived or deduced from without, but must nonetheless be assumed. It may well be the case, in other words, that Urobuchi is entirely content with his life, and writes about depressed characters and hopeless situations because he enjoys writing about them and knows he is good at it. We cannot know.

That said, Madoka Magica is quite easily readable as a story about depression. The magical girls all show different symptoms of depressive disorders, such as Mami's loneliness or Homura's lack of affect, and Sayaka's entire arc is a story of loss and mounting despair culminating in suicide. 

Depression and loss are everywhere in this first "true" episode of the series. Much of the first half of the episode concerns Madoka's attempts to deal with the loss of Mami. At breakfast the morning after Mami's death, the yolk of her egg reminds her of Mami's hair, and the taste brings her to tears in a scene strongly reminiscent of one of my favorite lines from another well-loved magical girl series:

"But I don't understand! I don't understand how this all happens. How we go through this. I knew her, and then she's—there's just a body, and I don't understand why she just can't get back in it and not be dead... anymore! It's stupid! It's mortal and stupid! And... I was having fruit punch, and I thought, well, [she] will never have any more fruit punch, ever, and she'll never have eggs, or yawn, or brush her hair, not ever, and no one will explain to me why!" \cite{ref39}

But Mami's loss and Madoka's mourning of her are hardly the only examples in the episode. Equally powerful is Kyousuke's loss of his ability, the grief and rage he expresses at the music he can only listen to but no longer play. In response to the revelation that his doctors have advised him to accept that his disability is permanent, Sayaka chooses to become a magical girl to save him, even after seeing how truly dangerous it is. For all that she will reveal mixed motives later, in this moment she is selfless and heroic. Kyousuke's inability to feel the injury he inflicts on his arm, unfortunately, is foreshadowing; Sayaka will soon likewise discover a state of numb despair that locks out pain at the price of locking out everything else with it. Additionally, Sayaka's choice to become Mami's replacement—that is, barreling ahead to save others without first addressing her own grief and loss—foreshadows that same behavior.

And of course there is the suicide cult. We see little of Hitomi throughout this series, so it is not clear just how sudden this is, but the implication seems to be that the witch drew together vulnerable people into this suicide pact. What exactly made Hitomi vulnerable is never explicitly answered, any more than it is for the unnamed woman in Episode 2. There are hints, however: prior to this episode Hitomi is characterized as being always busy, not only with the exam prep classes expected of a high-achieving student but also Japanese upper-class cultural training such as traditional dance or tea ceremony classes. There is no further reference to these classes after this episode, suggesting perhaps that she was overstressed and lightened her burden after this incident. Another possibility is that it involved shyness regarding her attraction to Kyousuke, which was apparently unknown even to her friends prior to her ultimatum to Sayaka shortly after this episode, following which she immediately approached Kyousuke.

More focus is on Madoka's vulnerability, as the witch preys on her failure to help Mami. Given the man in the suicide cult who talks about his failure to run a factory, and the possibilities for Hitomi's vulnerability, it seems that this witch's modus operandi is to target feelings of failure and inadequacy.

Its choice of method, at least against Madoka, is to confront her with her inadequacy by means of presenting her with television sets playing clips of the prior episode. As the witch does so, Madoka blurs and distorts, the usual outline that demarcates the character-background distinction having evaporated. In other words, the witch confronts Madoka with the fact that she is a television character, stripping away her identity. She is not a person, but a thing, one element in a television show, and thus begins dissolving away into the background of that show.

But it is a curious truth that, for all their fictionality, the feelings engendered by fictional people are nonetheless real. Even knowing that Madoka is fictional, we still care about her and are happy when Sayaka arrives to rescue her. For now, at least, this is enough to restore her to "reality" within the show. Madoka is still a distinct entity, not a diffuse concept, at least within the confines of the show. Nonetheless, the first episode of the "true" Madoka Magica foreshadows its ending; Madoka will be diffused.

By the close of the episode, the next arc is set fully into motion: Sayaka is a magical girl, her costume suggestive of a heroic knight protector, as she sets out to save Kyousuke, Madoka, and the world from evil witches. Unfortunately, arrayed against her is her perfect foil, experienced where she is naïve, red to her blue, self-centered and greedy to Sayaka's self-sacrificing protector. Or so, at least, the two of them appear to be; as we will see, like so many pairs of opposites in this show, they are not so different after all.
